% No preamble. Prefix csgq-. Fixed joint homogeneous background system.
\section{Summable connecting gauges and the invariant quotient}
\label{csgq-section}

This section closes the gauge-invariance interface for classical
substitution on the analytic interaction quotient. Its finite-regulator
inputs are the actual ordinary-minimum comparison and equality theorem,
the ordinary-to-corrected transfer, the constraint-preserving flat-slice
flow, and the coherent pre-gauge support estimates. It does not assume a
weighted estimate for the nonlinear gauge-fixing inverse.

\subsection{Joint limits, domains and the statement}
\label{csgq-input-section}

Fix a preceding homogeneous refinement $M$, a finite allowed endpoint
ratio $L$, and a single joint subsequence of compatible cubic regulators.
Let
\begin{equation}
 e^{ia_{L,N}(w)}=P_{M,L,N}(e^{iw})
       \longrightarrow e^{ia_L(w)}
 \label{csgq-background-limit}
\end{equation}
locally at every current link, locally uniformly in the finitely many
variables of each compactly supported final field $w$. The current
lattice has spacing $1/L$ in final physical units. Assume the established
coherent support bound
\begin{equation}
 \|a_L(w)|_i\|_\infty
 \le \rho e^{-\kappa\operatorname{dist}(i,K)},\qquad
 \rho=\frac{C_A b}{L},
 \qquad \|w\|_\infty<b,\quad \operatorname{supp}w\subset K.
 \label{csgq-background-decay}
\end{equation}
Here $K$ is a finite set of final cells. The same assertion holds for
every other compact datum in the ball. Joint analytic compactness and
the finite-regulator support estimates supply this input; unrelated
pointwise choices of subsequences do not replace it.

Use a unitarily invariant matrix norm for group elements. Its comparison
with the scalar-coordinate maximum norm defining the activity polydisc
has fixed color-dependent constants. In particular, constant conjugation
is not silently treated as an isometry of that coordinate maximum norm.
All such comparison constants are included below. The dimension is
$d\ge2$, with $d=4$ in the application. Constants in summability estimates
may depend on fixed $L$ and on a finite compact support, but not on
preceding refinement or approximating volume.

Let $\mathcal Q_{\rm inv}$ be the quotient of anchored analytic
interactions by equality of their vacuum-normalized totals on compact
fields. Its invariant totals are invariant under constant global
adjoint transformations and under compact local gauge paths that stay
inside the analytic field ball. An identity for arbitrary disconnected
endpoint gauges is not part of the conclusion. The summable-field
extension of the quotient gives, for $0<\rho_1<b$,
\begin{equation}
 |\Phi_E(A)-\Phi_E(B)|
 \le \frac{\|[E]\|_{\mathcal Q}}{b-\rho_1}
                 \|A-B\|_{\ell^1_{\rm cell}},
 \qquad \|A\|_\infty,\|B\|_\infty\le\rho_1.
 \label{csgq-summable-continuity}
\end{equation}
The normalized total is absolutely convergent on these summable fields,
and its value depends only on the class.

\begin{theorem}[Universal local gauge identity for a fixed joint map]
\label{csgq-main}
Under these inputs, the scalar substitution
\begin{equation}
 [E]\longmapsto
       \bigl(w\longmapsto \Phi_E(a_L(w))\bigr)
 \label{csgq-substitution}
\end{equation}
preserves local gauge invariance for every invariant input class on a
common smaller germ. If
\begin{equation}
             C_g C_A/L<1/2,
 \label{csgq-margin-condition}
\end{equation}
where $C_g$ is the fixed group/chart constant constructed below, the
compact connecting paths have a common strict margin in the original
radius-$b$ field ball. The local analytic identity then holds throughout
the connected common coarse domain:
\begin{equation}
             \Phi_E(a_L(w^v))=\Phi_E(a_L(w))
 \label{csgq-gauge-identity}
\end{equation}
for the allowed compact local gauge transformations.

The connecting gauges and domain choices are independent of $E$.
Consequently, when the separately established connected representation
and boundedness estimates for classical substitution apply,
\eqref{csgq-substitution} is an action on the whole invariant quotient
for this fixed joint background system.
\end{theorem}

\subsection{The small gauge supplied by the actual equality proof}
\label{csgq-finite-section}

We record the quantitative consequence of the ordinary-minimum equality
proof. For an ordinary-feasible fine field $U=e^{i\eta A}$, the first
transfer uses $g_1=e^{iH_b\theta}$. Its induced-gauge implicit equation
and the cumulative gauge $b_j(U)$ obey
\begin{equation}
 \|\theta\|_\infty\le2\|\omega\|_\infty,\qquad
 \omega=\log b_j(U)/i,\qquad
 \|\omega\|_\infty\le C\|A\|_{C^{1,\alpha}}.
 \label{csgq-first-transfer}
\end{equation}
The proved bound $H_b:\ell^\infty\to C^{2,\alpha}$ controls this gauge
uniformly in mesh and volume. The constraint-preserving flat-slice flow
has scalar velocity
\begin{equation}
 \lambda(t)=\mathcal L_{A(t)}A(t),\qquad
 \|\lambda(t)\|_{C^{2,\alpha}}
 \le C e^{-t}\|R_0D^T A(0)\|_{C^\alpha}
 \le C e^{-t}\|A(0)\|_{C^{1,\alpha}}.
 \label{csgq-flow}
\end{equation}
Integrating its actual unitary gauge ODE therefore gives
\begin{equation}
 \|g_2-I\|_\infty
 \le C\int_0^\infty\|\lambda(t)\|_\infty\,dt
 \le C\|A(0)\|_{C^{1,\alpha}}.
 \label{csgq-integrated-gauge}
\end{equation}
Finally the selected ordinary representative is obtained from the
strict corrected minimizer by
$h_*=e^{iE_{\rm pt}\omega_*}$, where
$\|\omega_*\|_\infty\le C\|B\|_\infty$ and
$E_{\rm pt}:\ell^\infty\to C^{2,\alpha}$ is the fixed smooth
point-value extension.

If equality holds in the two minimum comparisons, strict convexity on
the fixed slice identifies the intervening field with that corrected
minimizer. The product of the actual transfer gauges consequently maps
$U$ to the selected ordinary minimizer and satisfies
\begin{equation}
 \|g_{\rm eq}-I\|_\infty
       \le C\bigl(\|A\|_{C^{1,\alpha}}+\|B\|_\infty\bigr).
 \label{csgq-equality-gauge}
\end{equation}
The constants are uniform on the established smaller comparison balls.
This is stronger than equality of action values alone. It uses no
inverse massless gauge Laplacian and no unproved spatial locality of
$\mathcal L_A$.

For a compact coarse gauge $v=e^{i\theta}$, extend its prescribed values
to the original fine lattice by
$\widehat v=e^{iE_{\rm pt}\theta}$. This extension is compact in physical
space and has $C^{2,\alpha}$ norm bounded by $C\|\theta\|_\infty$,
uniformly in the fine mesh. The fields
$U_{\rm ord}(w)^{\widehat v}$ and $U_{\rm ord}(w^v)$ obey the same
ordinary constraint. The two minimum comparisons, one using the
inverse gauge, give equality. Equation \eqref{csgq-equality-gauge}
therefore provides an actual fine gauge $g_N$ with
\begin{equation}
 U_{\rm ord,N}(w^v)=U_{\rm ord,N}(w)^{g_N},\qquad
 \|g_N-I\|_\infty\le C(b+\|\theta\|_\infty)\le\delta.
 \label{csgq-finite-orbit}
\end{equation}
Choose the data and local gauge neighborhoods so that $\delta$ is a
fixed small number. Every competitor has complete fine-potential norm
at most $C(b+\|\theta\|_\infty)$ and hence lies in the actual comparison
class. Ordinary covariance of the prefix gives
\begin{equation}
 e^{ia_{L,N}(w^v)}
       =\bigl(e^{ia_{L,N}(w)}\bigr)^{g_N|_{\rm current}}.
 \label{csgq-prefix-orbit}
\end{equation}

\subsection{An infinite connecting gauge and its tail}
\label{csgq-tail-section}

Fix real compact $w,\theta$ in these neighborhoods. The endpoint
configurations already converge along the specified joint subsequence.
Compactness of the group and a diagonal subsequence over current sites
give a limiting gauge for this pair. Taking a further subsequence of
gauges changes neither specified endpoint limit. With
$U=e^{ia_L(w)}$ and $V=e^{ia_L(w^v)}$, local passage gives
\begin{equation}
 V_{xy}=g(x)U_{xy}g(y)^{-1},\qquad
                 \sup_x\|g(x)-I\|\le\delta.
 \label{csgq-infinite-orbit}
\end{equation}
No simultaneous or holomorphic choice of $g$ for all pairs is required.
Both endpoints have exponential decay because $w,w^v$ are compact.

\begin{lemma}[One constant at infinity]
\label{csgq-constant-lemma}
Let compact-group link fields $U,V$ on $\mathbb Z^d$, $d\ge2$, satisfy
$V=U^g$. Suppose their positive bonds based at $x$ obey
\begin{equation}
       \|U_b-I\|+\|V_b-I\|\le C_K e^{-\mu|x|_\infty}.
 \label{csgq-endpoint-tail}
\end{equation}
Then there is a single $c\in G$ such that
\begin{equation}
 \|g(x)-c\|\le C'_K(1+|x|_\infty)e^{-\mu|x|_\infty}.
 \label{csgq-gauge-tail}
\end{equation}
Thus $c^{-1}g-I$ is summable. If $\|g-I\|_\infty\le\delta$, its
supremum is at most $2\delta$.
\end{lemma}
\begin{proof}
The edge equation implies $g(y)=V_{xy}^{-1}g(x)U_{xy}$.
Unitary invariance of the matrix norm gives
\begin{equation}
       \|g(y)-g(x)\|\le\|U_{xy}-I\|+\|V_{xy}-I\|.
 \label{csgq-edge-bound}
\end{equation}
The ray $g(ne_1)$ is Cauchy, with tail $O(e^{-\mu n})$, and
its limit $c$ belongs to $G$. For $n=|x|_\infty$, the lattice boundary
of $[-n,n]^d$ connects $x$ to $ne_1$ by a path of length at most
$C_dn$. A boundary coordinate can be held at $\pm n$ while the other
coordinates are adjusted; a second boundary coordinate permits changing
its sign. This uses $d\ge2$. Each positive lower endpoint has norm at
least $n-1$. Summing \eqref{csgq-edge-bound} along the path gives
$C_dnC_Ke^{-\mu(n-1)}$. Combining this with the ray tail proves
\eqref{csgq-gauge-tail}. Polynomial shell growth proves summability.
Taking the limit in the small-supremum bound and using
$\|c^{-1}g-I\|=\|g-c\|$ gives the last assertion.
\end{proof}

For the actual fields, \eqref{csgq-background-decay} implies
\eqref{csgq-endpoint-tail} with $\mu$ a positive multiple of
$\kappa/L$, after enclosing their finite final supports in a cube.
Support-dependent constants are harmless for equality of evaluations.
In one dimension the two ends can have different constants; the
restriction $d\ge2$ is substantive.

\subsection{Compact gauge paths with a common strict margin}
\label{csgq-path-section}

Put $h=c^{-1}g$. Then $h-I\in\ell^1$ and
\begin{equation}
                         U^h=c^{-1}Vc.
 \label{csgq-normalized-gauge}
\end{equation}
For sufficiently small $\delta$ choose the common Hermitian logarithm
$h(x)=e^{i\lambda(x)}$. It has small supremum and belongs to $\ell^1$.
Let $\lambda_n$ equal $\lambda$ inside a current cube of radius $n$
and vanish outside, and set $h_n=e^{i\lambda_n}$. These gauges are
compact and
\begin{equation}
 \|h_n-h\|_{\ell^1}\longrightarrow0,\qquad
 \|U^{h_n}_b-U^h_b\|
       \le\|h_n(x)-h(x)\|+\|h_n(y)-h(y)\|.
 \label{csgq-compact-gauge-limit}
\end{equation}
Hence the transformed links, and their potentials on a common strict
logarithm chart, converge in cell $\ell^1$. A smooth physical cutoff is
not needed on this fixed current lattice: the truncation jump itself
decays exponentially.

The paths, not just their endpoints, stay in the analytic domain.
The edge equation and the Lipschitz bound for the small logarithm give
\begin{equation}
 \|\lambda(y)-\lambda(x)\|
 \le C\bigl(\|a_L(w)_b\|+\|a_L(w^v)_b\|\bigr)\le C\rho.
 \label{csgq-log-gradient}
\end{equation}
For all sufficiently large $n$, the exterior truncation jump of
$\lambda_n$ is also at most $\rho$. Duhamel's formula for real
Hermitian matrices yields, for $0\le t\le1$,
\[
 \|e^{it\lambda_n(y)}-e^{it\lambda_n(x)}\|
                  \le\|\lambda_n(y)-\lambda_n(x)\|.
\]
The small link logarithm and fixed color-norm comparisons consequently
give
\begin{equation}
 \sup_{0\le t\le1}
   \left\|\frac1i\log\bigl(U^{e^{it\lambda_n}}\bigr)\right\|_\infty
                    \le C_g\rho.
 \label{csgq-path-margin}
\end{equation}
Here $C_g$ is independent of support, truncation radius, refinement and
volume. Condition \eqref{csgq-margin-condition} places these paths
strictly below $b/2$. At a smaller fixed endpoint ratio, restricting
the data radius gives a local germ instead; a whole-ball margin is
not inferred without the condition.

The input field in \eqref{csgq-path-margin} is summable rather than
compact. Before invoking the defining invariance of the quotient,
truncate its potential to a much larger cube containing the support of
$\lambda_n$ and all adjacent bonds. Inside the gauge support this
compact input is unchanged. Outside, the gauge is the identity and
the field is either unchanged or zero. The same strict path bound
therefore holds. Compact-field local gauge invariance and then
\eqref{csgq-summable-continuity}, as the input truncation tends to
infinity, give
\begin{equation}
              \Phi_E(\log U^{h_n}/i)=\Phi_E(\log U/i).
 \label{csgq-compact-path-identity}
\end{equation}
Now send $n$ to infinity using \eqref{csgq-compact-gauge-limit} and
\eqref{csgq-summable-continuity}. Extend constant global adjoint
invariance to summable fields by the same compact-input approximation.
The result is
\begin{equation}
 \Phi_E(a_L(w))
   =\Phi_E(\operatorname{Ad}_{c^{-1}}a_L(w^v))
   =\Phi_E(a_L(w^v)).
 \label{csgq-real-identity}
\end{equation}
This proof uses the same gauges and margins for every invariant class.

\subsection{Analytic continuation and precise scope}
\label{csgq-scope-section}

To finish Theorem \ref{csgq-main}, differentiate
\eqref{csgq-real-identity} at zero compact coarse gauge parameter.
The resulting infinitesimal identity is holomorphic in the finitely
many coarse field variables on their connected product ball and
vanishes on a real open set. It therefore vanishes throughout that
complex ball. Integrating it along any allowed compact local gauge
path that stays in the ball proves \eqref{csgq-gauge-identity}.
This argument does not extend an infinitesimal Ward identity to
arbitrary disconnected endpoint transformations.

The same endpoint method preserves the spatial symmetries required
by the invariant quotient, when the actual finite blocking and
minimum problem have those symmetries. The transformed ordinary
minimizer and the selected minimizer at transformed data are small
feasible fields with equal minimum values. The actual equality proof
supplies the controlled gauge; prefix covariance gives the endpoint
orbit relation; compact data give exponential tails. The preceding
compact-gauge argument applies, and the input total's spatial symmetry
supplies the remaining change of variables. Constant adjoint covariance
is exact for the equivariant construction, or follows from the small
comparisons and connectedness of the compact group. Symmetry of
individual carrier coefficients is unnecessary.

Universality means that every admissible interaction class can be
evaluated using this one fixed joint infinite background system.
It does not assert uniqueness among different regulator subsequences.
Semigroup identities require the maps and their compositions to
belong to the same joint system; finite-regulator identities alone
do not establish consistency for unrelated limits.

The connected representation and boundedness of the quotient action
remain separate analytic inputs. This section proves their missing
universal gauge identity from the actual finite transfer and flow,
endpoint decay, and summable-field continuity. It concerns homogeneous
small fields, $d\ge2$, and compact local gauge paths. It proves neither
quantum-integral descent on the quotient nor quantum induction,
large-orbit or mixed-history control, or a continuum measure.
